\subsection*{1. Hough Transform and Template Matching}

Hough transforms just aren't built for finding Waldo. This method is looking for exact parameters, stuff like lines ($\rho, \theta$) or circles ($x, y, r$), and Waldo doesn't really have a parametric shape. He's basically just a random pattern of clothes and accessories. Even if you tried to match circles for his glasses, the scene is so crowded that the accumulator would end up with a ton of false peaks. And on top of that, Hough voting ignores color completely and takes forever to run on a big cluttered image.

Template matching isn't much better honestly. It assumes the target is the exact same scale and rotation as the template, so it falls apart the second there's a lighting change or something's occluded. The Waldo image also has way too many repeating patterns that end up scoring higher than the actual match. Basically this whole thing is a lot harder than the synthetic images from Part 1, since those objects were just sitting there isolated with nothing around them.

Out of the three metrics, \texttt{tm\_ssd} and \texttt{tm\_nssd} worked the best, they nailed the traffic light, the construction sign, and Waldo every time. Makes sense since they're just minimizing raw pixel differences, and Waldo's colors are pretty distinct from everyone around him. \texttt{tm\_nccor} on the other hand was all over the place. It drifted off the construction sign and totally whiffed on Waldo, locking onto some random red and black soldiers instead. I think the normalization is the problem, it makes the metric care more about local contrast than actual color, which kills it in a scene this repetitive.
